\section{Part 2 --- HIV Vaccination on the sIC AIDS Model}

This part takes the simplified Imperial College (sIC) AIDS compartment model from
Haefner~(2005), Chapter~15, and extends it with an HIV vaccination programme. The
task is essentially Exercise~15.6~\#5: vaccinate sexually active susceptibles into
a protected state at rate $\nu$, let protection wane at rate $l$, and then ask four
questions --- whether vaccination makes the epidemic peak and decline, what it
costs, what the optimum vaccination rate is, and how vaccination compares with
condom promotion. I work from the textbook equations and Table~15.2, the
vaccination design loosely follows Garnett \& Stover~(2002), and all numbers below
come from my own simulation of the model. Two of the modelling decisions needed to
make the textbook model behave as the textbook describes are not cosmetic, so I
state them explicitly and justify them rather than hide them.

\subsection{The sIC model}

The sIC model partitions the population into twelve compartments: the three disease
states susceptible ($S$), HIV-infected but pre-AIDS ($I$), and clinical AIDS ($A$),
crossed with two sexes (female $f$, male $m$) and two age classes (age~1 = birth to
15, pre-sexual; age~2 = 16 to death, sexually active). Only age-2 individuals
participate in transmission. The state vector is therefore
\[
\{S_{f,1},S_{f,2},S_{m,1},S_{m,2},\;
  I_{f,1},I_{f,2},I_{m,1},I_{m,2},\;
  A_{f,1},A_{f,2},A_{m,1},A_{m,2}\}.
\]
The model is a deliberate simplification of the full IC model: two sexes with
identical behaviour, a single AIDS stage, a single sexual-activity class (so the
social-mixing probability collapses to $\rho=1$), and two age classes.

\paragraph{Force of infection.}
HIV is sexually transmitted, so transmission is \emph{frequency-dependent}: people
seek a roughly fixed number of partners regardless of how large the population is,
so what matters is the \emph{proportion} of potential partners who are infectious,
not their absolute density. For a susceptible age-2 female the per-capita hazard is
(Eq.~15.6)
\begin{equation}
\lambda_{S_{f,2}} = c\,\rho\,\beta_{m,f}\,\frac{I_{m,2}}{S_{m,2}+I_{m,2}},
\qquad
\lambda_{S_{m,2}} = c\,\rho\,\beta_{f,m}\,\frac{I_{f,2}}{S_{f,2}+I_{f,2}},
\label{eq:foi}
\end{equation}
the second being the male analogue (Eq.~15.8). Here $c$ is the new-partner
acquisition rate, $\beta_{m,f}$ and $\beta_{f,m}$ are the per-partnership
transmission probabilities, and the factor $I/(S+I)$ is the frequency of
infectious individuals among possible partners. Two features deserve note. First,
the transmission probabilities are \emph{asymmetric}: $\beta_{m,f}=0.2$ versus
$\beta_{f,m}=0.075$, so male\,$\to$\,female transmission is about $2.7\times$ more
likely, a standard biological feature of heterosexual HIV. This is why women reach
higher equilibrium prevalence than men in the baseline run. Second, clinical-AIDS
individuals ($A$) are \emph{excluded} from the partner pool (assumed not sexually
active), so the denominator is $S+I$ only --- a point that turns out to matter for
the choice of epidemic seed.

\paragraph{The female equations.}
The full female subsystem (Eqs.~15.5a--f) is
\begin{align}
\dd{S_{f,1}}{t} &= \eta\theta\zeta\bigl(S_{f,2}+(1-\vartheta)I_{f,2}\bigr)
                   - \mu S_{f,1} - \xi S_{f,1}, \label{eq:Sf1}\\
\dd{S_{f,2}}{t} &= \xi S_{f,1} - (\lambda_{S_{f,2}}+\mu)\,S_{f,2}, \label{eq:Sf2}\\
\dd{I_{f,1}}{t} &= \eta\theta\zeta\vartheta\,I_{f,2}
                   - (\mu+\xi)I_{f,1} - \gamma I_{f,1}, \label{eq:If1}\\
\dd{I_{f,2}}{t} &= \lambda_{S_{f,2}} S_{f,2}
                   - (\mu+\gamma)I_{f,2} + \xi I_{f,1}, \label{eq:If2}\\
\dd{A_{f,1}}{t} &= \gamma I_{f,1} - (\mu+\xi+\alpha)A_{f,1}, \label{eq:Af1}\\
\dd{A_{f,2}}{t} &= \gamma I_{f,2} + \xi A_{f,1} - (\mu+\alpha)A_{f,2}. \label{eq:Af2}
\end{align}
The male equations~(15.7a--f) are identical in form, with the newborn sex split
$(1-\eta)$ replacing $\eta$ and the male force of infection
$\lambda_{S_{m,2}}$ replacing $\lambda_{S_{f,2}}$; note that male births are still
driven by the age-2 female compartments, since the model assumes males do not limit
reproduction.

The mechanisms encoded here are: \emph{reproduction}, with births created only from
age-2 females at fecundity $\theta$ (Eq.~\ref{eq:Sf1}); \emph{perinatal
transmission}, where an infected mother bears an infected newborn with probability
$\vartheta=0.35$ (feeding $I_{\cdot,1}$, Eq.~\ref{eq:If1}) and an uninfected
newborn with probability $1-\vartheta$ (feeding $S_{\cdot,1}$); \emph{ageing} at
rate $\xi$, moving each state from age~1 to age~2 within sex; \emph{AIDS
progression} at rate $\gamma$ moving $I\to A$; \emph{AIDS mortality} $\alpha$ added
to natural mortality $\mu$ in the $A$ compartments; and \emph{demography}, with
natural death $\mu$ acting on every compartment. The parameters are listed in
Table~\ref{tab:sic-params}.

\begin{table}[t]
\centering
\caption{Parameters of the simple IC (sIC) AIDS model, after Haefner~(2005),
Table~15.2 (values attributed there to Garnett \& Anderson, 1993). Variable
``values'' are the initial conditions used for the baseline run.}
\label{tab:sic-params}
\begin{tabular}{llrl}
\toprule
Symbol & Meaning & {Value} & Unit \\
\midrule
$\alpha$        & AIDS-stage extra death rate                 & 1.0    & yr$^{-1}$ \\
$\beta_{f,m}$   & female$\to$male transmission probability    & 0.075  & --- \\
$\beta_{m,f}$   & male$\to$female transmission probability     & 0.2    & --- \\
$c$             & new-partner acquisition rate                 & 2.35   & yr$^{-1}$ \\
$\eta$          & proportion of newborns female                & 0.5    & --- \\
$\gamma$        & progression rate $I\to A$ (Table value)      & 1.16   & yr$^{-1}$ \\
$\mu$           & natural (non-AIDS) death rate                & 0.0227 & yr$^{-1}$ \\
$\rho$          & social-mixing probability                    & 1.0    & --- \\
$\theta$        & female fecundity                             & 0.2088 & yr$^{-1}$ \\
$\vartheta$     & perinatal transmission probability           & 0.35   & --- \\
$\xi$           & ageing rate age-1$\to$age-2                   & 0.0667 & yr$^{-1}$ \\
$\zeta$         & proportion in sexual-activity class          & 1.0    & --- \\
\midrule
\multicolumn{4}{l}{\emph{Initial conditions (individuals):}
  $S_{f,1}=S_{m,1}=3000$, $S_{f,2}=S_{m,2}=1000$,} \\
\multicolumn{4}{l}{\quad all $I=0$, all $A=0$ except the seed (see
  \S\ref{sec:judgments}).} \\
\bottomrule
\end{tabular}
\end{table}

\paragraph{A conservation correction to the printed equations.}
The AIDS age-2 equation deserves a note. The textbook prints the ageing inflow to
$A_{f,2}$ (and $A_{m,2}$) as $\xi I_{f,1}$, but this violates conservation of
individuals. In Eq.~\ref{eq:Af1}, $A_{f,1}$ \emph{loses} $\xi A_{f,1}$ to ageing,
and that flux must land in $A_{f,2}$; meanwhile the ageing of $I_{f,1}$ is
\emph{already} counted as an inflow to $I_{f,2}$ in Eq.~\ref{eq:If2}. Using
$\xi I_{f,1}$ in Eq.~\ref{eq:Af2} would therefore double-count the $I_{f,1}$ cohort
and lose the $A_{f,1}$ cohort. I use the conservative form $\xi A_{f,1}$
(and $\xi A_{m,1}$) shown above. As a check, the total-population balance
$\frac{\mathrm{d}}{\mathrm{d}t}\sum(\text{states}) = \text{births}-\text{deaths}$
then holds to $\sim\!10^{-13}$ in the simulation, which would not be the case with
the printed term.

\subsection{Two modelling judgments}
\label{sec:judgments}

Two decisions were forced on me by trying to make the model reproduce the
textbook's stated behaviour. I present both honestly, because each changes the
results.

\paragraph{(i) The progression rate $\gamma$.}
For a two-sex frequency-dependent model the basic reproduction number is the
geometric mean of the two cross-sex reproduction numbers,
\begin{equation}
\Rzero = \frac{c\,\sqrt{\beta_{m,f}\,\beta_{f,m}}}{\mu+\gamma}.
\label{eq:R0}
\end{equation}
The Table~15.2 value $\gamma=1.16\,\text{yr}^{-1}$ implies a mean time from HIV
infection to AIDS of $1/(\mu+\gamma)\approx0.85$~yr --- less than a year, which
contradicts the textbook's own statement that progression takes 1--10~years. Worse,
substituting it into Eq.~\ref{eq:R0} gives $\Rzero\approx0.24<1$, so the epidemic
cannot ignite at all; a literal-parameter simulation confirms the seeded infection
dies out (peak prevalence $\sim\!7\times10^{-4}$). This directly contradicts
Fig.~15.5, where HIV persists endemically. I therefore adopt the biologically
grounded $\gamma=0.1\,\text{yr}^{-1}$, giving a mean infectious period of
$1/(\mu+\gamma)\approx8.1$~yr --- comfortably inside the textbook's stated 1--10~yr
range --- and
\begin{equation}
\Rzero = \frac{2.35\,\sqrt{0.2\times0.075}}{0.0227+0.1} \approx 2.35 > 1,
\end{equation}
so the epidemic invades, as the textbook figure requires. The associated
herd-immunity (critical vaccinated) fraction is $p_c = 1-1/\Rzero \approx 0.574$.

\paragraph{(ii) The epidemic seed.}
Table~15.2 seeds the epidemic with five AIDS males ($A_{m,2}=5$). But AIDS
individuals are excluded from the force of infection (Eq.~\ref{eq:foi}), so a seed
placed entirely in $A$ cannot transmit anything: the five seed males simply
progress to death without infecting anyone, and the epidemic never starts (verified
numerically --- maximum incidence exactly zero). I therefore seed five
\emph{infectious} males ($I_{m,2}=5$) instead. This is the smallest change that lets
the documented epidemic ignite while respecting the model's own rule that $A$ does
not transmit.

\subsection{Validation against Fig.~15.5}

With the two corrections above, the baseline ($\nu=0$) run reproduces the
qualitative behaviour of Fig.~15.5a (\verb|verify_baseline_prevalence.png|): the
total population grows over the early horizon (from $N(0)=8005$ to a peak of
$\sim\!22{,}900$ around year~48, driven by the high developing-country fertility),
HIV does \emph{not} burn out but settles to an endemic level, and female prevalence
exceeds male prevalence. The simulated equilibrium prevalences are
\begin{center}
female $\approx 0.77$, \quad male $\approx 0.61$, \quad overall $\approx 0.41$,
\end{center}
against the textbook targets of $\sim\!0.9$ (female) and $\sim\!0.75$ (male). The
ordering and endemic persistence match; the absolute levels run somewhat lower,
which is unsurprising given that I had to recalibrate $\gamma$ to make the epidemic
ignite at all. The time to exceed 10\% overall prevalence is about 34~yr, in the
same decades-long ballpark as the textbook's $\sim\!20$~yr.

\begin{figure}[t]
\centering
\includegraphics[width=0.78\linewidth]{verify_baseline_prevalence.png}
\caption{Baseline sIC validation ($\nu=0$). HIV persists endemically and female
prevalence exceeds male prevalence, reproducing the qualitative behaviour of
Fig.~15.5a; equilibrium prevalences are female $\approx0.77$, male $\approx0.61$.}
\label{fig:verify-baseline}
\end{figure}

The condom scenario (Fig.~15.5b) is modelled, as in the textbook, by halving both
transmission probabilities ($\beta_{f,m}:0.075\to0.0375$,
$\beta_{m,f}:0.2\to0.1$). This multiplies $\Rzero$ by exactly $0.5$ (it enters
Eq.~\ref{eq:R0} as $\sqrt{\beta_{m,f}\beta_{f,m}}$), giving $\Rzero\approx1.17$ ---
only just above one. In my calibration the resulting equilibrium prevalence is very
low ($\sim\!0.001$, \verb|verify_condom_prevalence.png|), which is a larger drop
than the textbook's Fig.~15.5b ($\sim\!0.7$ female, $\sim\!0.45$ male). The
direction is right (condoms lower prevalence and lengthen the time-to-10\%) but the
magnitude diverges, and the reason is instructive: at $\gamma=0.1$ the halved-$\beta$
$\Rzero$ sits barely above the threshold $\Rzero=1$, so the epidemic is in the
extremely sensitive near-threshold regime where small changes in calibration move
the endemic level a great deal. The textbook's larger residual prevalence is
consistent with its (implicitly) higher underlying $\Rzero$. I flag this divergence
rather than tune it away; it is the same near-threshold sensitivity that recurs in
Q(d).

\begin{figure}[t]
\centering
\includegraphics[width=0.78\linewidth]{verify_condom_prevalence.png}
\caption{Condom scenario (both $\beta$ halved). Prevalence collapses because
$\Rzero$ drops to $\approx1.17$, just above threshold. The direction matches
Fig.~15.5b; the magnitude is larger because the halved-$\beta$ system sits in the
sensitive near-threshold regime (see text).}
\label{fig:verify-condom}
\end{figure}

\subsection{The vaccination extension}

The vaccine adds two protected compartments, $P_{f,2}$ and $P_{m,2}$, at age~2 only
(age~2 is the terminal age class, so there is no ageing out of $P$). Susceptibles
are vaccinated at per-capita rate $\nu$ and protection wanes back to susceptible at
rate $l$:
\begin{align}
\dd{P_{f,2}}{t} &= \nu\,S_{f,2} - (l+\mu)\,P_{f,2}, &
\dd{P_{m,2}}{t} &= \nu\,S_{m,2} - (l+\mu)\,P_{m,2}, \\
\dd{S_{f,2}}{t} &= \xi S_{f,1} - (\lambda_{S_{f,2}}+\mu+\nu)\,S_{f,2} + l\,P_{f,2}, &
\dd{S_{m,2}}{t} &= \xi S_{m,1} - (\lambda_{S_{m,2}}+\mu+\nu)\,S_{m,2} + l\,P_{m,2},
\end{align}
with all other equations unchanged. A cost accumulator tracks the cumulative number
of vaccinations administered,
\begin{equation}
\dd{V}{t} = \nu\bigl(S_{f,2}+S_{m,2}\bigr),
\qquad \text{cost} = \$10\times V,
\label{eq:cost}
\end{equation}
integrated alongside the state so there is no post-hoc quadrature error. Each \$10
charge counts a vaccination \emph{event}, which correctly includes re-vaccinating
people whose protection waned and returned them to $S$. With $\nu=0$ the protected
and cost compartments stay identically zero and the model reduces exactly to the
baseline.

\paragraph{Vaccine type.}
I model a \emph{take} (all-or-nothing) vaccine with waning: while in $P$ an
individual is treated as fully protected (no force of infection acts on $P$, and
$P$ never flows to $I$), and a fraction loses protection at rate $l$. This is the
cleanest reading of the requested $S\to P\to S$ flow and maps onto the IC model's
``fully immunised + waning'' category. The project values $\nu=0.65\,\text{yr}^{-1}$
and $l=0.1\,\text{yr}^{-1}$ correspond to the Garnett/Stover ``65\% coverage'' and
``10-year mean duration'' standard scenario; using $\nu$ as a constant rate rather
than a coverage level reached after five years is a deliberate simplification.
Ignoring demography, the protected fraction saturates at a \emph{waning ceiling}
\begin{equation}
\frac{P}{S+P} \;\to\; \frac{\nu}{\nu+l} = \frac{0.65}{0.75} \approx 0.867,
\end{equation}
so even at this rate the vaccine can never protect everyone --- at most about 87\%
of the eligible pool at any instant.

\paragraph{The partner-pool convention.}
One modelling choice in the extension is genuinely load-bearing: whether the
protected individuals belong in the denominator of the frequency-dependent force of
infection. Protected people are sexually active and uninfected, so they \emph{are}
legitimate partners --- a partnership with a protected person simply transmits
nothing. The correct frequency-dependent denominator is therefore
$S+I+P$ (with $A$ still excluded). I adopt this convention. It matters because
moving susceptibles into $P$ then \emph{dilutes} the infectious frequency
$I/(S+I+P)$, which is exactly the mechanism of herd immunity. If instead one
excluded $P$ from the denominator, vaccinees would be removed from \emph{both} the
numerator and the denominator, artificially holding the infectious fraction up and
spuriously erasing the herd-immunity threshold --- vaccination would then appear to
have no critical rate at any level. That is an artifact of the denominator choice,
not a property of the disease; I report it only as a sensitivity check in Q(c). At
$\nu=0$ this choice is moot ($P=0$), so the baseline and condom scenarios are
unaffected.

\subsection{Q(a): Does vaccination cause a peak and decline?}

The question asks whether vaccination makes the epidemic peak and then decline. The
answer, at $\nu=0.65\,\text{yr}^{-1}$, is actually stronger than ``peak and
decline'': because the achievable protected fraction (0.867) comfortably exceeds the
herd-immunity threshold $p_c=0.574$, vaccination pushes the effective reproduction
number below one,
\begin{equation}
\Reff \approx \Rzero\Bigl(1-\tfrac{\nu}{\nu+l}\Bigr)
       \approx 2.35\times(1-0.867) \approx 0.31 < 1,
\end{equation}
so the seeded epidemic \emph{never ignites} in the first place.

The contrast with the baseline is stark (Table~\ref{tab:qa},
Figs.~\ref{fig:qa-incidence}--\ref{fig:qa-protected}). The baseline incidence rises
to a large peak of $\sim\!791$ new infections/yr around year~48 and stays high,
whereas under vaccination incidence is highest at $t=0$ (just the five seed males
infecting a handful of contacts) and then \emph{falls monotonically} to essentially
zero. Prevalence shows the same picture: the baseline climbs to a $\sim\!0.41$
endemic plateau, while the vaccinated run shows only a negligible interior bump
($\sim\!0.001$ around year~4.6) before declining to near zero. The decline is not
instantaneous because the existing $I$/$A$ stock carries the virus for the $\sim\!8$--$10$~yr
it takes to progress and die; but with $\Reff<1$ there is no fresh transmission to
refill it, so the infection clears. Meanwhile the protected fraction of the eligible
age-2 pool rises and saturates at $\approx0.85$, just below the analytic ceiling
$0.867$ (slightly lower because natural mortality $\mu$ also acts on $P$),
confirming the take-with-waning steady state.

\begin{table}[t]
\centering
\caption{Q(a): epidemic outcome with and without vaccination ($\nu=0.65$,
80-yr horizon).}
\label{tab:qa}
\begin{tabular}{lcc}
\toprule
 & Baseline ($\nu=0$) & Vaccine ($\nu=0.65$) \\
\midrule
Peak incidence (new/yr)      & 790.6 (@ 47.8 yr) & 2.3 (@ $t=0$, the seed) \\
Incidence behaviour          & rise to peak, stay high & monotone decline \\
Peak overall prevalence      & 0.409 & 0.0010 (@ $\approx$4.6 yr) \\
Prevalence @ 80 yr           & 0.408 & $\approx 0$ \\
Protected fraction $P/(S+P)$ & --- & $\to 0.85$ \\
\bottomrule
\end{tabular}
\end{table}

\begin{figure}[t]
\centering
\includegraphics[width=0.78\linewidth]{qa_incidence.png}
\caption{Q(a): incidence (new infections/yr), vaccine versus baseline. The baseline
peaks at $\sim\!791$/yr near year~48; under vaccination ($\Reff\approx0.31$)
incidence declines monotonically from the seed and never takes off.}
\label{fig:qa-incidence}
\end{figure}

\begin{figure}[t]
\centering
\includegraphics[width=0.78\linewidth]{qa_prevalence.png}
\caption{Q(a): overall prevalence, vaccine versus baseline. The baseline rises to a
$\sim\!0.41$ endemic plateau; the vaccinated run shows only a tiny early bump before
declining to near zero.}
\label{fig:qa-prevalence}
\end{figure}

\begin{figure}[t]
\centering
\includegraphics[width=0.78\linewidth]{qa_protected_fraction.png}
\caption{Q(a): protected fraction of the eligible age-2 pool, saturating near the
analytic waning ceiling $\nu/(\nu+l)=0.867$.}
\label{fig:qa-protected}
\end{figure}

\subsection{Q(b): Cost of vaccination}

Using the cost accumulator of Eq.~\ref{eq:cost}, Table~\ref{tab:qb} reports the
cumulative vaccinations, total cost, infections averted (cumulative baseline
infections minus cumulative vaccinated infections), and cost per infection averted
at several horizons.

\begin{table}[t]
\centering
\caption{Q(b): cost and cost-effectiveness of vaccination ($\nu=0.65$, \$10/vax).}
\label{tab:qb}
\begin{tabular}{rrrrr}
\toprule
$T$ (yr) & Vaccinations & Cost (\$) & Infect.\ averted & \$/infection averted \\
\midrule
20 & 11{,}374 & 113{,}738 &    305 & 372 \\
30 & 16{,}217 & 162{,}165 &  1{,}585 & 102 \\
50 & 24{,}746 & 247{,}461 & 13{,}434 &  18 \\
\bottomrule
\end{tabular}
\end{table}

Total cost grows nearly linearly: $\sim\!\$114$k at 20~yr, $\sim\!\$162$k at 30~yr,
$\sim\!\$247$k at 50~yr (Fig.~\ref{fig:qb}). Most of the expense is the up-front
vaccination of the standing susceptible pool; once the protected fraction
saturates, the programme keeps re-vaccinating waned individuals and new entrants at
a roughly steady $\sim\!407$ vaccinations/yr, i.e.\ a recurring cost of only
$\sim\!\$4{,}070$/yr.

Cost-effectiveness improves sharply over time. Early on the untreated baseline has
barely taken off, so there are few infections to avert and the ratio looks expensive
(\$372 per infection averted at 20~yr). As the baseline epidemic accelerates toward
its $\sim\!791$/yr incidence peak while the vaccinated arm stays near zero, averted
infections grow far faster than cost, and the cost per infection averted falls to
\$102 at 30~yr and \$18 at 50~yr. The 30-yr figure of \$102 sits just below the
\$110--390 band that Stover \& Garnett~(2002) report for the Imperial-College /
rural-Zimbabwe setting (on their \$20/person basis), which is consistent given the
\$10 unit cost here and the fact that, under the corrected partner-pool convention,
vaccination averts nearly the entire baseline epidemic rather than only part of it.
The absolute dollar totals depend on this small synthetic population, so the
transferable metric is the cost \emph{per infection averted}, not the headline cost.

\begin{figure}[t]
\centering
\includegraphics[width=0.78\linewidth]{qb_cost.png}
\caption{Q(b): (top) cumulative cost and vaccinations versus time, with horizon
markers at 20/30/50~yr; (bottom) cumulative infections for baseline versus vaccine
and the infections-averted band.}
\label{fig:qb}
\end{figure}

\subsection{Q(c): Optimum vaccination rate}

``Optimum'' has two defensible meanings, and I report both, because they give
different answers.

\paragraph{Epidemiological threshold $\nu_c$.}
The smallest rate that drives $\Reff<1$ follows from requiring the protected
fraction to reach $p_c$. Since the steady protected fraction is $\nu/(\nu+l)$, the
condition $\nu/(\nu+l)\ge p_c$ gives a critical rate
\begin{equation}
\nu_c = \frac{l\,p_c}{1-p_c},
\qquad p_c = 1-\tfrac{1}{\Rzero}\approx0.574,
\end{equation}
which evaluates to $\nu_c\approx0.135\,\text{yr}^{-1}$. This closed form is the
right order of magnitude and direction but is a lower bound, because it omits two
mechanisms the full sIC model contains: perinatal transmission ($\vartheta=0.35$),
an infection route vaccination does not block, and demographic recruitment with
mortality acting on $P$ (which holds the realised protected fraction slightly below
the ceiling). Measuring the threshold directly --- seeding a tiny infection into the
vaccinated disease-free population and finding the $\nu$ at which the early growth
rate crosses zero ($\Reff=1$) --- gives a simulated $\nu_c\approx0.366\,\text{yr}^{-1}$
(Fig.~\ref{fig:qc-threshold}). The standard $\nu=0.65$ sits comfortably above this,
so it genuinely controls the epidemic. I emphasise that this is a frequency-dependent
\emph{prevalence-proportion} threshold rather than an ``$(I+A)/N\to0$'' criterion:
the severe AIDS mortality ($\alpha=1\,\text{yr}^{-1}$) drives the whole population
toward extinction over centuries even at $\nu=0$, so absolute counts collapse for
the wrong reason, and the invasion / $\Reff$ test is the robust signal.

\paragraph{Cost-effective rate.}
Sweeping $\nu$ over a 50-yr horizon, infections averted \emph{saturate} past
$\nu_c$: below it each extra unit of $\nu$ averts many infections, but once
$\nu\gtrsim\nu_c$ essentially all of the achievable $\sim\!13{,}400$ infections are
already averted and further vaccination only re-vaccinates waned individuals for
negligible benefit (Figs.~\ref{fig:qc-prev},~\ref{fig:qc-ce}). The
diminishing-returns knee --- the smallest $\nu$ capturing 99\% of the maximum
achievable aversion --- is at $\nu\approx0.175\,\text{yr}^{-1}$, with a marginal cost
there of $\sim\!\$122$ per extra infection averted. The standard $\nu=0.65$ still
controls the epidemic (it is above $\nu_c$) but over-vaccinates somewhat relative to
this cost-effective knee; its average cost is \$18 per infection averted at 50~yr.

\paragraph{Sensitivity.}
The exact threshold tracks $\Rzero$: raising the partner rate to $c=3.0$ lifts
$\Rzero$ to $\sim\!3.0$ and the simulated $\nu_c$ to $\sim\!0.48$; lowering it via
$\gamma=0.18$ drops $\Rzero$ to $\sim\!1.42$ and $\nu_c$ to $\sim\!0.16$. The
conclusion is robust in direction --- a finite $\nu_c=l\,p_c/(1-p_c)$ always exists
under the correct convention and $\nu=0.65$ stays above it across this plausible
range --- but near-threshold cases (small $\Rzero$) leave the margin thin.

\paragraph{Caveat.}
If $P$ is wrongly excluded from the partner pool, the invasion growth rate never
crosses zero in the sweep: vaccination would appear to have \emph{no} finite
threshold at any rate and merely shrink the epidemic. As argued above this is an
artifact of dropping uninfected, sexually-active vaccinees out of the
frequency-dependent denominator, and I report it only to show that the threshold
result is genuine, not an accident of the denominator choice.

\begin{figure}[t]
\centering
\includegraphics[width=0.80\linewidth]{qc_invasion_threshold.png}
\caption{Q(c): invasion growth rate $r(\nu)$ crossing zero ($\Reff=1$) under the
corrected convention ($P$ in pool, finite $\nu_c\approx0.366$); the wrong
$P$-excluded convention (no crossing) is shown for contrast. Theory $\nu_c$,
simulated $\nu_c$, and $\nu=0.65$ are marked.}
\label{fig:qc-threshold}
\end{figure}

\begin{figure}[t]
\centering
\includegraphics[width=0.78\linewidth]{qc_prevalence_vs_nu.png}
\caption{Q(c): peak epidemic size versus $\nu$. Above the threshold the epidemic is
suppressed.}
\label{fig:qc-prev}
\end{figure}

\begin{figure}[t]
\centering
\includegraphics[width=0.78\linewidth]{qc_cost_effectiveness.png}
\caption{Q(c): average and marginal (incremental) cost per infection averted versus
$\nu$, with the diminishing-returns knee at $\nu\approx0.175$ marked.}
\label{fig:qc-ce}
\end{figure}

\subsection{Q(d): Vaccination versus condoms}

Finally, which intervention controls the epidemic better: vaccination ($\nu=0.65$,
$l=0.1$) or safe-sex education and condoms (modelled, as in the textbook, by halving
both transmission probabilities)? They act on different levers --- condoms lower
$\Rzero$ directly through $\beta$, vaccination moves people into $P$ --- so I compare
them on matched outcomes (Table~\ref{tab:qd}, Figs.~\ref{fig:qd-comp},~\ref{fig:qd-averted}).

\begin{table}[t]
\centering
\caption{Q(d): strategy comparison (80-yr horizon, $\Rzero=2.35$ baseline).}
\label{tab:qd}
\begin{tabular}{lcrr}
\toprule
Strategy & $\Rzero$ / $\Reff$ & Infect.\ averted & Cost (\$) \\
\midrule
Baseline                       & $\Rzero=2.35$         & ---       & 0 \\
Vaccination ($\nu=0.65$)       & $\Reff=0.31$          & 31{,}224  & 348{,}618 \\
Condoms ($0.5\beta$)           & $\Rzero=1.17$         & 30{,}857  & 0 (unpriced) \\
\bottomrule
\end{tabular}
\end{table}

Both strategies essentially control the epidemic. Vaccination drives
$\Reff=0.31<1$ (invasion rate $-0.040\,\text{yr}^{-1}$): it strictly crosses the
elimination threshold because the protected fraction $0.867$ exceeds $p_c=0.574$.
Condoms put $\Rzero$ at $1.17$ --- still marginally \emph{above} one (invasion rate
$+0.086\,\text{yr}^{-1}$), so a 50\% $\beta$ cut is not quite enough to cross the
threshold at this $\Rzero$, though the epidemic is so slow that it is effectively
controlled within the horizon. On infections averted the two are nearly identical:
vaccination averts $31{,}224$ of the baseline $31{,}235$, condoms $30{,}857$, with
vaccination edging ahead only because it pushes $\Reff$ strictly below one (residual
$11$ cumulative infections versus $378$ for condoms). Both collapse peak prevalence
from $0.41$ to $\sim\!0.001$ with no large interior peak.

Where they differ sharply is cost basis, and this is where I decline to crown a
single winner. Vaccination carries an \emph{explicit} \$348{,}618 over 80~yr and is
capped by the waning ceiling $0.867$, so it must re-vaccinate waned individuals
indefinitely. Condoms are modelled as a permanent $0.5\times$ cut to $\Rzero$ at no
priced cost --- but that cost is \emph{unpriced}, not zero; a real behaviour-change
programme has its own budget the project does not supply. A dollar-matched
comparison would need a condom-programme cost we do not have. There is also the
near-threshold caveat: condoms leave $\Rzero$ only $\sim\!17\%$ above one, and a
higher baseline $\Rzero$ (say a higher partner rate) would push condom-$\Rzero$
clearly above one while raising vaccination's $\nu_c$ --- both verdicts are
$\Rzero$-sensitive near threshold, which is exactly why I report $\Reff$ explicitly
rather than just trajectories. The honest reading is a near-tie on epidemiological
outcome, with vaccination marginally ahead for strictly crossing the threshold, and
the two not directly comparable on cost as priced here.

\begin{figure}[t]
\centering
\includegraphics[width=0.78\linewidth]{qd_strategy_comparison.png}
\caption{Q(d): overall prevalence trajectories for baseline, vaccination, and
condoms. Both interventions collapse the epidemic; the baseline rises to a
$\sim\!0.41$ plateau.}
\label{fig:qd-comp}
\end{figure}

\begin{figure}[t]
\centering
\includegraphics[width=0.78\linewidth]{qd_averted_and_reff.png}
\caption{Q(d): infections averted and reproduction number per strategy. Vaccination
crosses $\Reff<1$ strictly; condoms sit just above $\Rzero=1$.}
\label{fig:qd-averted}
\end{figure}

\subsection{Part 2 summary}

On the sIC model, with $\gamma$ recalibrated to a biologically sensible
$0.1\,\text{yr}^{-1}$ (giving $\Rzero\approx2.35$) and the epidemic seeded in the
infectious rather than the AIDS compartment, vaccination at $\nu=0.65\,\text{yr}^{-1}$
controls the HIV epidemic: it drives $\Reff\approx0.31<1$, so the epidemic never
ignites, and prevalence stays near zero instead of rising to the $\sim\!0.41$
endemic plateau of the baseline. It is cost-effective ($\sim\!\$102$ per infection
averted at 30~yr, within the published Garnett/Stover range), though the standard
rate over-vaccinates somewhat relative to the cost-effective knee at
$\nu\approx0.18\,\text{yr}^{-1}$. Against condom promotion the two are roughly
comparable on epidemiological outcome --- both essentially control the epidemic,
vaccination crossing the threshold strictly while condoms hover just above it --- but
they are not directly comparable on cost as priced here, and both verdicts are
sensitive to the baseline $\Rzero$ near threshold. All of these conclusions rest on
the documented modelling assumptions, in particular the $\gamma$ recalibration and
the partner-pool convention that gives vaccination its herd-immunity effect.
